% p334_1; source PDF page 334; manually verified editable table.
\begin{tikzpicture}[x=1cm,y=-1cm,inner sep=0pt,outer sep=0pt]
\path[use as bounding box] (0,0) rectangle (16.000,9.445);
\draw[line width=.4pt] (0.000,0.000) rectangle (2.450,1.215);
\node[align=center,text width=2.190cm,font={\linespread{1}\fontsize{10}{13.50}\selectfont\bfseries}] at (1.225,0.608) {食品犯罪};
\draw[line width=.4pt] (2.450,0.000) rectangle (4.850,1.215);
\node[align=center,text width=2.140cm,font={\linespread{1}\fontsize{10}{13.50}\selectfont\bfseries}] at (3.650,0.608) {成立条件};
\draw[line width=.4pt] (4.850,0.000) rectangle (7.500,1.215);
\node[align=center,text width=2.390cm,font={\linespread{1}\fontsize{10}{13.50}\selectfont\bfseries}] at (6.175,0.608) {食品类型};
\draw[line width=.4pt] (7.500,0.000) rectangle (16.000,1.215);
\node[align=center,text width=8.240cm,font={\linespread{1}\fontsize{10}{13.50}\selectfont\bfseries}] at (11.750,0.608) {范例};
\draw[line width=.4pt] (0.000,1.215) rectangle (2.450,6.030);
\node[align=center,text width=2.190cm,font={\linespread{1}\fontsize{10}{13.50}\selectfont}] at (1.225,3.623) {生产、销售有毒、有害食品罪};
\draw[line width=.4pt] (2.450,1.215) rectangle (4.850,6.030);
\node[align=center,text width=2.140cm,font={\linespread{1}\fontsize{10}{13.50}\selectfont}] at (3.650,3.623) {抽象危险犯（也即行为犯）};
\draw[line width=.4pt] (4.850,1.215) rectangle (7.500,6.030);
\node[align=left,text width=2.390cm,font={\linespread{1}\fontsize{10}{13.50}\selectfont}] at (6.175,3.623) {有毒、有害食品，是指掺入有毒、有害的非食品原料的食品};
\draw[line width=.4pt] (7.500,1.215) rectangle (16.000,6.030);
\node[align=left,text width=8.240cm,font={\linespread{1}\fontsize{10}{13.50}\selectfont}] at (11.750,3.623) {1.用“瘦肉精”喂猪或销售这样的猪肉。注意：“食品”包括食材，如生猪、活鱼、生长中的蘑菇等。（2022年试题）\\{}2.将工业酒精勾兑成散装白酒出售。（2022年试题）\\{}3.将工业用猪油冒充食用油出售。\\{}4.用“地沟油”生产食用油};
\draw[line width=.4pt] (0.000,6.030) rectangle (2.450,9.445);
\node[align=center,text width=2.190cm,font={\linespread{1}\fontsize{10}{13.50}\selectfont}] at (1.225,7.738) {生产、销售不符合安全标准的食品罪};
\draw[line width=.4pt] (2.450,6.030) rectangle (4.850,9.445);
\node[align=center,text width=2.140cm,font={\linespread{1}\fontsize{10}{13.50}\selectfont}] at (3.650,7.738) {具体危险犯};
\draw[line width=.4pt] (4.850,6.030) rectangle (7.500,9.445);
\node[align=left,text width=2.390cm,font={\linespread{1}\fontsize{10}{13.50}\selectfont}] at (6.175,7.738) {不符合安全标准的食品，没有达到毒害性};
\draw[line width=.4pt] (7.500,6.030) rectangle (16.000,9.445);
\node[align=left,text width=8.240cm,font={\linespread{1}\fontsize{10}{13.50}\selectfont}] at (11.750,7.738) {1.含有严重超出标准限量的农药残留、重金属、污染物质的食品。\\{}2.包装材料、容器被污染的食品。\\{}3.超过保质期的食品};
\end{tikzpicture}
